% r03_r1_geo: the cumulant bound for the step Gaussian, part III: the four shifted pavements on Z^3, on the cubic tori and on the skew tori, and the layers of a cell.

\begin{lemma}[pavements]\label{rot:lem:R1geo}
Let $b\ge b_{\mathrm{geo}}\defeq576$, $q_b\defeq\lceil b^2\rceil$, $w_b\defeq\lceil b^{3/2}\rceil$ and $s_b\defeq\lfloor q_b/4\rfloor$, so that $s_b\ge5w_b+2$, $q_b\ge6w_b$ and $3q_b\le4b^2$; put $v_b\defeq(4b^2)^3=64\,b^6$. On each $\mathbb{L}$ of the setting there are four partitions $\mathcal P^i$, $i=0,1,2,3$, of $\mathbb{L}$ into \emph{cells} $\Box$, with \emph{levels} $\lambda_i:\mathbb{L}\to[0,\infty]$, such that
\begin{enumerate}[label=\textup{(G\arabic*)},leftmargin=*]
\item $\lambda_i$ is $1$-Lipschitz;
\item if $x,y$ lie in different cells of $\mathcal P^i$, then $|x-y|\ge\lambda_i(x)+\lambda_i(y)$;
\item $\bigcap_{i=0}^3\{x:\ \lambda_i(x)<2.5w_b+1\}=\emptyset$;
\item every cell has at most $v_b$ sites; in particular, if some $\mathcal P^i$ consists of a single cell, then $|\mathbb{L}|\le v_b$.
\end{enumerate}
\end{lemma}
\begin{proof}
\emph{Numbers.} For $b\ge576$: $s_b\ge b^2/4-1\ge5b^{3/2}+7\ge5w_b+2$ (the middle inequality is $\sqrt b\ge20+32b^{-3/2}$, true for $\sqrt b\ge24$), $q_b\ge b^2\ge6b^{3/2}+6\ge6w_b$, and $3q_b\le3b^2+3\le4b^2$.

\emph{Coordinates.} Each $\mathbb{L}$ has three coordinate cycles $\mathbb I_a$, $a=1,2,3$, with a well-defined map $x\mapsto x_a\in\mathbb I_a$ and $N_a\defeq|\mathbb I_a|\in\N\cup\{\infty\}$: for $\Z^3$, $\mathbb I_a=\Z$; for $(\Z/N\Z)^3$, $\mathbb I_a=\Z/N\Z$; for $\Z^3/mG_\pm$, $\mathbb I_a\defeq\Z/N_a\Z$ with $N_a\Z\defeq\pi_a(mG_\pm)$ the projection of the period lattice on the coordinate $a$, that is $N_1=N_2=10m$ and $N_3=50m$ (the entries $40,\mp30,0$ and $\pm30,40,0$ have greatest common divisor $10$); the map is well defined because the period lattice $G\defeq mG_\pm$ is contained in $G'\defeq\bigoplus_aN_a\Z$. In every case $\dist_{\mathbb I_a}(x_a,y_a)\le|\tilde x_a-\tilde y_a|\le|\tilde x-\tilde y|=|x-y|$ for lifts realising the quotient distance. For the skew tori we record: $\det G=(40\cdot40+30\cdot30)\cdot50\cdot m^3=125000\,m^3=(50m)^3$, $[G':G]=125000/5000=25$, and $G\cap\{g_3=0\}=m\langle(40,\mp30,0),(\pm30,40,0)\rangle$ has index $2500m^2$ in $\Z^2\times\{0\}$.

\emph{Cuts.} A coordinate $a$ is \emph{cut} if $N_a\ge q_b$ (always for $\Z^3$). For a cut coordinate choose a set $B^0_a\subset\mathbb I_a+\frac12$ of cuts all of whose (cyclic) gaps lie in $[q_b,2q_b)$: for $\mathbb I_a=\Z$ take $q_b\Z+\frac12$; on a cycle take $p\defeq\lfloor N_a/q_b\rfloor\ge1$ cuts with gaps $\lfloor N_a/p\rfloor$ or $\lceil N_a/p\rceil$, which lie in $[q_b,2q_b-1]$. Put $B^i_a\defeq B^0_a+is_b$, $i=0,\dots,3$. Cuts of $B^i_a$ and of $B^{i'}_a$, $i\ne i'$, are at cyclic distance at least $s_b$: two shifts of the same cut differ by $|i-i'|s_b\in[s_b,3s_b]$, and $N_a-3s_b\ge q_b-3s_b\ge s_b$; shifts of different cuts differ by at least $q_b-3s_b\ge s_b$. For $x\in\mathbb{L}$ put $\lambda_i(x)\defeq\min_{a\ \mathrm{cut}}\dist_{\mathbb I_a}(x_a,B^i_a)$, and $\lambda_i\defeq+\infty$ if no coordinate is cut. The \emph{cut set} $C^i$ is the set of points $\xi$ of the continuum torus ($\R^3$, $\R^3/N\Z^3$ or $\R^3/G$) with $\xi_a\in B^i_a$ for some cut coordinate $a$, and the cells of $\mathcal P^i$ are the intersections of $\mathbb{L}$ with the connected components of the complement of $C^i$; no site lies on a cut, so the cells partition $\mathbb{L}$.

\emph{(G1)} is the Lipschitz property of $\dist_{\mathbb I_a}(\cdot,B^i_a)$ together with $\dist_{\mathbb I_a}(x_a,y_a)\le|x-y|$. \emph{(G3)}: if $\lambda_i(x)<2.5w_b+1$ for all four $i$, each minimum is attained at a cut coordinate, and since there are at most three, two indices $i\ne i'$ use the same coordinate $a$; then a cut of $B^i_a$ and a cut of $B^{i'}_a$ are at cyclic distance less than $2(2.5w_b+1)=5w_b+2\le s_b$, a contradiction. If no coordinate is cut the four sets are empty. \emph{(G2)}: lift $x,y$ to $\tilde x,\tilde y\in\R^3$ with $|\tilde x-\tilde y|=|x-y|$. The lift $\tilde C^i$ of the cut set is the union of the planes $\{\xi_a=c\}$ over the cut coordinates $a$ and the lifts $c$ of the cuts of $B^i_a$; it is $G'$-periodic, hence $G$-periodic, the components of $\R^3\setminus\tilde C^i$ are the products of open gap intervals in the cut coordinates and of $\R$ in the others, and the projection maps them onto the components of the complement of $C^i$. If the segment $[\tilde x,\tilde y]$ met none of the planes it would lie in one component, and $x,y$ would lie in one cell; so it crosses a plane $\{\xi_a=c\}$ with $c$ between $\tilde x_a$ and $\tilde y_a$, and $|x-y|\ge|\tilde x_a-\tilde y_a|=|\tilde x_a-c|+|c-\tilde y_a|\ge\dist_{\mathbb I_a}(x_a,B^i_a)+\dist_{\mathbb I_a}(y_a,B^i_a)\ge\lambda_i(x)+\lambda_i(y)$. \emph{(G4)}: a component $Q$ of $\R^3\setminus\tilde C^i$ projects onto its cell injectively modulo $G_Q\defeq\{g\in G:\ Q+g=Q\}=\{g\in G:\ g_a=0\text{ for every cut }a\}$: if $Q$ and $Q+g$ meet they coincide, $\tilde C^i$ being $G$-invariant. On $\Z^3$, and whenever all three coordinates are cut, $G_Q=0$ and the cell is a box with sides at most $2q_b-1$, hence has at most $(2q_b)^3$ sites. On a cubic torus the only other case is that no coordinate is cut, when $|\mathbb{L}|=N^3<q_b^3$. On $\Z^3/mG_\pm$, where $N_3=5N_1$, two further cases occur: only the third coordinate is cut ($10m<q_b\le50m$), the cell is the product of $\Z^2/(G\cap\{g_3=0\})$ with an interval of fewer than $2q_b$ integers and has fewer than $2500m^2\cdot2q_b=25\,(10m)^2\cdot2q_b<50\,q_b^3$ sites; or no coordinate is cut ($50m<q_b$) and $|\mathbb{L}|=\det G=(50m)^3<q_b^3$. With $q_b\le b^2+1\le\frac{257}{256}b^2$ the three counts are below $8.1\,b^6$, $50.6\,b^6$ and $1.02\,b^6$, hence below $v_b=64\,b^6$.
\end{proof}

\paragraph{Layers.} For a cell $\Box$ of $\mathcal P^i$ write $\lambda\defeq\lambda_i$ and define, all intersected with $\Box\cap\Lambda$:
\[
\Gamma_1(\Box)\defeq\{\lambda<w_b\},\quad\Gamma_2(\Box)\defeq\{w_b\le\lambda<2w_b\},\quad\Gamma_3(\Box)\defeq\{1.5w_b\le\lambda<2w_b\},\quad\Box'\defeq\{\lambda\ge2w_b\},\quad\Gamma_4(\Box)\defeq\{2w_b\le\lambda<2.5w_b\},
\]
$\mathrm{Deep}(\Box)\defeq\{\lambda\ge2.5w_b\}$ and $D_\Box\defeq\Gamma_2(\Box)\cup\Box'=\{\lambda\ge w_b\}$. By (G1) and (G2), two points of one cell at levels $\lambda(x)\ge\lambda(y)$ are at distance at least $\lambda(x)-\lambda(y)$, and two points of different cells at distance at least $\lambda(x)+\lambda(y)$. Hence: the sets $D_\Box$ of distinct cells are at distance at least $2w_b$; $\dist(D_\Box,\Gamma_1(\tilde\Box))\ge w_b$ and $\dist(\mathrm{Deep}(\Box),\tilde\Box\cap\{\lambda\ge1.5w_b\})\ge4w_b$ for $\tilde\Box\ne\Box$; and $\dist(\mathrm{Deep}(\Box),\Lambda\setminus\Box')\ge\frac12w_b$, $\dist(\Box\cap\{\lambda\ge1.5w_b\},\Lambda\setminus D_\Box)\ge\frac12w_b$, $\dist(\mathrm{Deep}(\Box),\Lambda\setminus(\Box\cap\{\lambda\ge1.5w_b\}))\ge w_b$, $\dist(\Gamma_4(\Box),\Box\cap\{\lambda<1.5w_b\})\ge\frac12w_b$, $\dist(\Box',\Gamma_2(\Box)\setminus\Gamma_3(\Box))\ge\frac12w_b$, $\dist(\Gamma_1(\Box),\Box')\ge w_b$.
